\section{The residual representation of a row and its upward theta image}

% BEGIN J86O1
\paragraph{J86O1. Finding the Hecke point at the same coefficient place.}
Fix an original member, a location $i$ and a nonzero row
$M_j[\sigma]$, with $j=j(i)$ and $\lambda=w(i)$.
Take its theta type $P$ from U86T4. Its image is nonzero and
lies in $V_P=\mathrm{typeV}(i,P)$. We need a point for this
lift at some location with the same pair $(j,\lambda)$.

Use the original $\mathrm{NdSpec}$ for this pair. Write $O_c$ for
its common coefficient DVR, $K_c$ for its fraction field, and $T_c$
for the $O_c$-algebra generated by the specified split spherical
operators on the full rational Hecke space.
By M84T7 each generator acts on $V_P$ by a scalar, the same for all
the allowed frames and uniformizers in its test. The rational basis
of forms gives both the complex basis and the $K_c$-basis by the
member's two basis identifications. In those bases each spherical
operator is the image of the same rational matrix, by
$h_L$ and $h_K$ in the Hecke model. The embedding of $K_c$ in
$\mathbb C$ is the extension of $\mathrm{Sp.iotaC}$, compatible
with those bases by $\mathrm{phK\_emb}$ and
$\mathrm{iotaC\_coh}$. Each such matrix splits over $K_c$ by the
full generalized eigenspace decomposition. Its eigenvalues belong to
$O_c$ by $\mathrm{Sp.std}$. In particular all scalars on $V_P$
are images of elements of $O_c$.

Choose one nonzero vector in $V_P$. The simultaneous equations for
these eigenvalues over $K_c$ have a nonzero solution there: their
linear functionals span a finite-dimensional dual subspace, a finite
basis of which has the same matrix rank after extension to
$\mathbb C$. Acting on a nonzero solution yields a character
$\psi:T_c\to O_c$. Indeed, each element is an $O_c$-polynomial
in the operators, hence acts by its evaluation at the integral
scalars. A relation among the operators kills this vector, so that
its scalar evaluation vanishes. The character fixes $O_c$ and its
simultaneous eigenspace is nonzero. Its residue kernel is maximal,
since the constants surject onto the residue field.
The member has $\mu_{\rm At}(\lambda)\ne\mu_0$ at every actual
location. Thus $\mathrm{Sp.mx\_all}$ gives a location $i'$ of
the same row and place with this maximal ideal, and
$\mathrm{Sp.ptX}$ gives $x'\in X_{i'}$ for $\psi$.

The tuple at both locations is $\mu_{\rm Tw}(\lambda)$; the
representation occurs in the same automorphic space with the same
level. By surjectivity onto its occurring irreducible types, choose
$P'$ at $i'$ representing this same type. The two type subspaces are
identified, as are the original theta operators on $M_j$. Thus
$P'$ is a theta type of $\sigma$. Its generators act by the
values of $x'$ by $\mathrm{Sp.ptX\_hT}$ and the common coefficient
embedding, on every vector of the type. The point test of M84T7
gives $\mathrm{ptV}(i',P')=x'$.
% END J86O1

% BEGIN J86O2
\paragraph{J86O2. Joint occurrence with the prescribed residue map.}
Take any $k,v,i,\sigma$ in the original joint occurrence test,
with $g=g_k\ne\ell_i$, and put $w=\mathrm{plGw}(k,v)$,
$\gamma=\mathrm{tg}(k,v)$. Use $i',P',x'$ from J86O1.
By U86T4 choose a row list $\Sigma$ at $i'$ whose augmentation
is a base-change list of $P'$. Goodness of $\sigma$ and the
nonzero row make it a singleton cusp list. At all designated
places the list has exactly the row labels, since these places
split and $\mathrm{lab\_loc}$ reads their labels. It follows that
$\Sigma$ satisfies both goodness tests with the original characters.
The distinguished term in its augmentation puts $x'$ in $Z$,
and the same list witnesses $x'\in Z_{\rm good}$. Apply R86G7 at
this point to obtain the \emph{untwisted} integral representation
$\rho$ of $\Sigma$ over $O_{i'}$.

The original datum has the identity $D.\mathrm{rbar\_tie}$:
for every list satisfying $\mathrm{HasParamW}$ and every integral
$\mathrm{IsRepOfSig}$ representation, it identifies the
characteristic polynomial of $D.\mathrm{ntypes.rbar}(\sigma,w(i'))$
at \emph{every} group element with the reduction of that integral
polynomial through $\mathrm{kFbar}_{i'}$. This applies to the pair
just constructed. Both $w(i')$ and $w(i)$ equal $\lambda$, so the
residual matrices in the required test are the very same matrices.
The identification of residual characteristics is the canonical one
from $w(i')=\lambda$, $\ell_{i'}=\operatorname{resChar}(\lambda)$.

For the reduction of the roots in R86G2 take the member's own
valuation ring $B=A.\mathrm{pl.plB}(\lambda)$ in $\mathbb C$ and
the residue identification $\kappa=A.\mathrm{pl.pl}\kappa(\lambda)$.
The equality $\mathrm{Au.kFbar}$ makes $\kappa$ after reduction
on $O_{i'}$ the specified map on integral coefficients.
The equality $A.\mathrm{iota}$ makes it on each Teichm\"uller
value $\tau(a)$ exactly $D.\mathrm{ntypes.iota}(k,v,\lambda)(a)$.
All these roots of unity are units of $B$. Reduction kills their
$\ell_i$-power factors. Thus with $L=D.\mathrm{nd.lab}(k,\sigma,v)$
and $a_{\ell_i'}$ the prime-to-$\ell_i$ part, R86G1--2 give
\[
 \det\!\left(X-\bar r_{\sigma,\lambda}(\gamma)\right)
   =\prod_{a\in\operatorname{roots}(L)}
        \left(X-D.\mathrm{ntypes.iota}(k,v,\lambda)(a_{\ell_i'})\right),
\]
including all multiplicities in the original field.

The wild inertia polynomial of $\rho$ is $(X-1)^r$ by the full
degree tame reading in R86G1--2. The same polynomial holds for the
actual $\bar r_{\sigma,\lambda}$ by $D.\mathrm{rbar\_tie}$. The
native residual representation has open kernel, hence its image of
wild inertia is a finite $g$-group. Each of its matrices is unipotent
in characteristic $\ell_i$ and has $\ell_i$-power order. The two
orders are coprime, so this image is trivial. Arithmetic Frobenius at
the designated place acts on tame inertia by the $g$th power since
the completion is $\mathbb Q_g$. A lift $\phi$ in the decomposition
group consequently satisfies
$\bar r_{\sigma,\lambda}(\phi\gamma\phi^{-1})
=\bar r_{\sigma,\lambda}(\gamma)^g$. This proves
$\mathrm{RowsC.c4b}$ on the original member.
% END J86O2

% BEGIN J86O3
\paragraph{J86O3. The upward inclusion at each specified point.}
Fix $i$, $x\in Z_{\rm good}$ and any $\Sigma$ satisfying
$\mathrm{HasParamX}(i,x,\Sigma)$. Choose a witness $Q$ with
point $x$ and with the prescribed augmented list. It suffices to
test a summand $M_j[\sigma]$ of $M_{j,\Sigma}$, so assume
$\mathrm{HasParamW}(i,\sigma,\Sigma)$ and the summand is nonzero.
By U86T4 its upward image lies in some $V_P$ and the lift $P$
has the augmented base change of a row list $\Sigma'$.
At each split place the two row lists $\Sigma,\Sigma'$ read
the same parameter $\mathrm{locW}(i,\sigma,w)$. Appending the
same distinguished string and tensoring by $\eta^{-1}$ give
the same local parameter, by the direct-sum and tensor-product
formulas. Hence $\mathrm{locV}(i,P,w)=\mathrm{locV}(i,Q,w)$
at each split place.

At each place in the split-good test the hyperspecial operator
of index $b$ acts by a scalar on $V_P$ by M84T7. The member's
spherical polynomial formula identifies $(-1)^b(\mathrm Nw)^{b(b-1)/2}$
times this scalar with the coefficient of $X^{n_V-b}$ of the
normalized local polynomial of $P$. Indeed one may evaluate on a
nonzero fixed vector in the irreducible representation of $P$:
project a nonzero level vector in $V_P$ onto one of its irreducible
copies and intertwine this copy with the chosen representative.
The projected vector remains hyperspecial fixed and has the original
Hecke scalars, since the operators are sums of translations. This
proves the coefficient identity for every frame and uniformizer in
the point test. The polynomial just identified equals the one for
$Q$, which is the Hecke polynomial of $x$ by $\mathrm{pt\_loc}$.
Each coefficient factor is nonzero, so all tested scalars on $V_P$
are the specified ones for $x$. The member's converse point test
gives $\mathrm{ptV}(i,P)=x$. Consequently the upward image of this
summand lies in $V_P\subseteq\mathrm{eigC}(i,x)$. Summing over
$\sigma$ proves $\mathrm{RowsC.parUp}$ for every tested $x,\Sigma$.
% END J86O3
